%%%%%%%%%%%%%%%%%%%%%%%%%%%%%%%%%%%%%%%%%%%%%%%%%%%%%%%%%%%%%%%%%%%%%
% Talk for visit of Philipps Research in Hamburg about work at KIT
% related to my dissertation. Approximately 15-20 min including
% discussion. 
%%%%%%%%%%%%%%%%%%%%%%%%%%%%%%%%%%%%%%%%%%%%%%%%%%%%%%%%%%%%%%%%%%%%%%
\documentclass{beamer}
%\usetheme{Warsaw}
%\usetheme{CambridgeUS}

\setbeamertemplate{navigation symbols}{} % don't use navigation tools on slides
%\setbeamertemplate{caption}{\raggedright\insertcaption\par}
%\usepackage[font=small,skip=1pt,labelformat=empty]{caption}
%%%%%%%%%%%%%%%%%%%%%%%%%%%%%%%%%%%%%%%%%%%%%%%%%%%%%%%%%%%%%%%%%%%%%%
%% PACKAGES
%%%%%%%%%%%%%%%%%%%%%%%%%%%%%%%%%%%%%%%%%%%%%%%%%%%%%%%%%%%%%%%%%%%%%%
\usepackage[utf8]{inputenc}
\usepackage[T1]{fontenc}
\usepackage[ngerman,swedish,british]{babel}
% from beamer example
\usepackage[absolute,overlay]{textpos}
\usepackage{extrabeamercmds}
 \usepackage{%
%  ae,
%  times,
%  arial,
  lmodern,
  charter,
%  multibib,
%  bibunits,
%  graphicx,
%  picins,
}
\usepackage{tikz}
\usetikzlibrary{shapes.arrows}
\usepackage{siunitx}
\PassOptionsToPackage{fleqn}{amsmath}% math environments and more by the AMS 
 \usepackage{amsmath}
\usepackage{verbatim}
%\usepackage{euler}
\usepackage[euler-digits]{eulervm}
\usefonttheme{professionalfonts}
% Convert eps to pdf
\PassOptionsToPackage{process=auto}{pstool}%[process=none]%auto,all
 \usepackage{pstool}
%%%%%%%%%%%%%%%%%%%%%%%%%%%%%%%%%%%%%%%%%%%%%%%%%%%%%%%%%%%%%%%%%%%%%%
%% COMMANDS
%%%%%%%%%%%%%%%%%%%%%%%%%%%%%%%%%%%%%%%%%%%%%%%%%%%%%%%%%%%%%%%%%%%%%%
%% From phdthesis
\newcommand{\ie}{i.\,e.}
\newcommand{\Ie}{I.\,e.}
\newcommand{\eg}{e.\,g.}
\newcommand{\Eg}{E.\,g.} 
% Use \boldsymbol instead of \mathbf to allow for bold greek letter.
% The package bm is more carefull, but changes font of bodl letter.
% Use \mathbold to use bodl fonts from eulermath
%\renewcommand{\vec}[1]{\mathbf{#1}}
%\renewcommand{\vec}{\boldsymbol}
\renewcommand{\vec}{\mathbold}
% Vectors
\newcommand{\xt}{{\vec{x},t}}
% 1d
\newcommand{\kp}{k_\perp}
\newcommand{\xip}{\xi_{\perp}}
\newcommand{\xipm}{\abs{\xi_{\perp}}_m}
\newcommand{\xin}{\xi_{\mathrm{noise}}}
\newcommand{\xim}{\xi_{m}}
\newcommand{\xic}{\xi_{\mathrm{c}}}
\newcommand{\xib}{\xi_{\mathrm{b}}}
\newcommand{\xir}{{\xi_{r}}}
\newcommand{\xis}{{\xi_{s}}}
\newcommand{\xix}{{\xi_{x}}}
\newcommand{\xiy}{{\xi_{y}}}
\newcommand{\xiz}{{\xi_{z}}}
% Vectors: 2D
\newcommand{\np}{\nabla_\perp}
\newcommand{\inp}{\nabla_\perp^{-2}}
\newcommand{\npa}{\nabla_{\perp,\alpha}}
\newcommand{\vecdp}{\vec{d}_\perp}
\newcommand{\veckp}{\vec{k}_\perp}
\newcommand{\vecqp}{\vec{q}_\perp}
\newcommand{\vecrp}{\vec{r}_\perp}
\newcommand{\vecxp}{\vec{x}_\perp}
\newcommand{\vecsp}{\vec{s}_\perp}
\newcommand{\vecxip}{\vec{\xi}_\perp}
% Vectors: 3D
\newcommand{\vecd}{\vec{d}}
\newcommand{\vecj}{\vec{j}}
\newcommand{\veck}{\vec{k}}
\newcommand{\vecn}{\vec{n}}
\newcommand{\vecp}{\vec{p}}
\newcommand{\vecq}{\vec{q}}
\newcommand{\vecr}{\vec{r}}
\newcommand{\vecs}{\vec{s}}
\newcommand{\vecv}{\vec{v}}
\newcommand{\vecx}{\vec{x}}
\newcommand{\x}{\vec{x}}
\newcommand{\vecxi}{\vec{\xi}}
\newcommand{\vecA}{\vec{A}}
\newcommand{\vecB}{\vec{B}}
\newcommand{\vecE}{\vec{E}}
\newcommand{\vecD}{\vec{D}}
\newcommand{\vecH}{\vec{H}}
\newcommand{\vecM}{\vec{M}}
\newcommand{\vecP}{\vec{P}}
\newcommand{\vecS}{\vec{S}}
% Other
\newcommand{\BP}{\operatorname{BP}}
\newcommand{\FBP}{\operatorname{FBP}}
\newcommand{\adu}{\operatorname{ADU}}
\newcommand{\dqe}{\operatorname{DQE}}
\newcommand{\sgn}{\operatorname{sgn}}
\newcommand{\snr}[1]{\operatorname{SNR_{#1}}}
\newcommand{\NA}{\mathrm{NA}}
\newcommand{\Nrings}{N_{\circledcirc}}
\newcommand{\Npix}{N_\mathrm{x}}
\newcommand{\fwhm}{\operatorname{FWHM}}
\newcommand{\tie}{\operatorname{\textcolor{blue}{TIE}}}
\newcommand{\ctf}{\operatorname{\textcolor{red}{CTF}}}
\newcommand{\qp}{\operatorname{\textcolor{green}{QP}}}
\newcommand{\qpc}{\operatorname{\textcolor{black}{QP_\mathrm{con}}}}
\newcommand{\conv}{\operatorname{\ast\ast}}
\newcommand{\Prop}{\mathcal{P}}
\newcommand{\order}{\mathcal{O}}
\newcommand{\N}{\mathcal{N}}
\newcommand{\R}{\mathbb{R}}
\newcommand{\sigb}{\sigma_\mathrm{b}}
\newcommand{\sigs}{\sigma_\mathrm{s}}
\newcommand{\sig}[1]{\sigma_\mathrm{#1}}
%\newcommand{\Seik}{\mathcal{S}}
\newcommand{\e}[1]{\mathrm{e}^{#1}}
\newcommand{\me}{m_e}
\newcommand{\ed}{\rho_{e}}
\newcommand{\Iin}{I^{\mathrm{(in)}}}
\newcommand{\Ibin}{\bar{I}^{\mathrm{(in)}}}
\newcommand{\Psin}{ \Psi^\mathrm{(in)} }
\newcommand{\subn}[1]{_\mathrm{#1}}
\newcommand{\submax}{_\mathrm{max}}
\newcommand{\supin}{^\mathrm{(in)}}
\newcommand{\supl}{^{(l)}}
\newcommand{\env}{^\mathrm{(env)}}
\newcommand{\ho}[1]{^{(#1)}}
\newcommand{\hon}[1]{^{(\mathrm{#1})}}
\newcommand{\vecjin}{\vec{j}^{\mathrm{(in)}}}
\newcommand{\vecjsc}{\vec{j}^{\mathrm{(s)}}}
\newcommand{\hook}{k_h}
\newcommand{\dpo}{\delta\phi_0}
\newcommand{\dpom}{\abs{\delta{\phi_0}}_{\mathrm{max}}}
\newcommand{\delo}{\delta_\omega}
\newcommand{\fov}{\mathrm{FOV}}
\newcommand{\lmaxg}{{l_{\mathrm{max},g}}}
\newcommand{\lmaxphi}{{l_{\mathrm{max},\phi}}}
\newcommand{\nlo}[1]{\text{\small{NLO}}_{#1}}
\newcommand{\lc}{l_\mathrm{c}}
%\newcommand{\textsub}[1]{\ensuremath{_{\mbox{\footnotesize #1}}}}
%\newcommand{\sqrd}{{\textsuperscript 2}}
\newcommand{\rar}{\rightarrow}
\newcommand{\lar}{\leftarrow}
\newcommand{\ra}{\right\rangle}
\newcommand{\la}{\left\langle}
\newcommand{\ii}{\mathrm{i}}
% modulus
\newcommand{\abs}[1]{\left| #1 \right|} % for absolute value
\newcommand{\absn}[1]{| #1 |} % for absolute value when | too big
% average 
\newcommand{\avg}[1]{\left\langle #1 \right\rangle} % for average
\newcommand{\avgn}[1]{\langle #1 \rangle} 
% mean abs
\newcommand{\avgabs}[1]{\left\langle\left|#1\right|\right\rangle} 
\newcommand{\avgabsn}[1]{\langle|#1|\rangle} 
\newcommand{\mabs}[1]{\overline{#1}} 
% Fourier operator and symbol
%\newcommand{\F}{\operatorname{\mathcal{F}}}
\newcommand{\F}{\mathcal{F}}
\newcommand{\iF}{\operatorname{\mathcal{F}^{-1}}}
\newcommand{\FO}[1]{\operatorname{\mathcal{F}#1}}
\newcommand{\ft}[1]{\widehat{#1}}
%\newcommand{\ft}[1]{\widetilde{#1}}
% angular averaged fft
\newcommand{\aft}[1]{\widetilde{#1}} 
% quasiparticle modification
\newcommand{\qpft}[1]{\widehat{#1}^{\mathrm{(QP)}}}
% derivatives and integral measures
\newcommand{\pd}{\partial}
\newcommand{\difd}{\mathrm{d}}
\newcommand{\Int}{\int\limits}
\newcommand{\intd}{\!\mathrm{d}}
\newcommand{\Intd}[1]{\!\mathrm{d}#1\,}
\newcommand{\Intdd}[1]{\!\mathrm{d^2}#1\,}
\newcommand{\Intddd}[1]{\!\mathrm{d^3}#1\,}
\newcommand{\const}{\mathrm{constant}}
% Matrices
\newcommand{\vv}[2]{\left(\begin{array}{c}#1\\#2\end{array}\right)}
\newcommand{\vvv}[3]{\left(\begin{array}{c}#1\\#2\\#3\end{array}\right)}
\newcommand{\mmmm}[4]{\left(\begin{array}{cc}#1&#2\\#3&#4\end{array}\right)}
% Other
\renewcommand{\hbar}{\hslash}
\newcommand{\bra}[1]{\left\langle#1\right|}
\newcommand{\ket}[1]{\left|#1\right\rangle}
% parentheses, brackets, and braces
\newcommand{\lp}{\left(}
\newcommand{\rp}{\right)}
\newcommand{\lrp}[1]{\left(#1\right)}
\newcommand{\lrb}[1]{\left[#1\right]}
\newcommand{\lrbr}[1]{\left\{#1\right\}}
\newcommand{\lb}{\left[}
\newcommand{\rb}{\right]}
\newcommand{\expsb}[1]{\exp\left[#1\right]}
\newcommand{\expb}[1]{\exp\left[#1\right]}
\newcommand{\expp}[1]{\exp\left(#1\right)}
%% BEAMER or FILE specific
\tikzset{
    myarrow/.style={
        draw,
        fill=black,
        single arrow,
        minimum height=3ex,
        single arrow head extend=1ex
    }
}
\newcommand{\arrowright}{%
  \tikz [baseline=-0.5ex]{\node [myarrow,rotate=0] {};} }
\newcommand{\sref}[1]{{\scriptsize [#1]}}
\newcommand{\srefr}[1]{\begin{flushright}{\tiny [#1]}\end{flushright}}
\newcommand{\flr}[1]{\begin{flushright}{#1}\end{flushright}}
%%%%%%%%%%%%%%%%%%%%%%%%%%%%%%%%%%%%%%%%%%%%%%%%%%%%%%%%%%%%%%%%%%%%%%
%% BEAMER CONFIG
%%%%%%%%%%%%%%%%%%%%%%%%%%%%%%%%%%%%%%%%%%%%%%%%%%%%%%%%%%%%%%%%%%%%%%
\newcommand{\logoheight}{0.05\textwidth}

\title[X-ray phase-contrast tomography]{In vivo imaging using \\X-ray
  phase-contrast microtomography} 

%\subtitle{Visit to Philips Hamburg} 

\author{Julian Moosmann}

%\institute{Laboratory for Applications of Synchrotron radiation (LAS)}

\date{August 17, 2015, Philips Hamburg}

\titlegraphic{\includegraphics[
%width=\textwidth,
height=.4\textheight
]{figures/cover/Nature}}

\addtobeamertemplate{title page}{\hfill 
  \includegraphics[height=\logoheight]{figures/logo/kth} \hfill 
  \includegraphics[height=\logoheight]{figures/logo/kit} \hfill 
  \includegraphics[height=\logoheight]{figures/logo/anka} \hfill 
  \includegraphics[height=\logoheight]{figures/logo/anl} \hfill 
  \includegraphics[height=\logoheight]{figures/logo/esrf} \hfill 
}{}

\AtBeginSubsection[]
{
  \begin{frame}[plain]{Outline}
    \tableofcontents[currentsection,currentsubsection]
  \end{frame}

} 

%%%%%%%%%%%%%%%%%%%%%%%%%%%%%%%%%%%%%%%%%%%%%%%%%%%%%%%%%%%%%%%%%%%%%%
%% BEGIN DOCUMENT
%%%%%%%%%%%%%%%%%%%%%%%%%%%%%%%%%%%%%%%%%%%%%%%%%%%%%%%%%%%%%%%%%%%%%%
\begin{document}
\selectlanguage{british}

%%%%%%%%%%%%%%%%%%%%%%%%%%%%%%%%%%%%%%%%%%%%%%%%%%%%%%%%%%%%%%%%%%%%%%
%% TITLE PAGE
\begin{frame}
  \titlepage
\end{frame}

%%%%%%%%%%%%%%%%%%%%%%%%%%%%%%%%%%%%%%%%%%%%%%%%%%%%%%%%%%%%%%%%%%%%%%
%% TABLE OF CONTENTS
\frame{
  \frametitle{Outline}
  \tableofcontents
}

%%%%%%%%%%%%%%%%%%%%%%%%%%%%%%%%%%%%%%%%%%%%%%%%%%%%%%%%%%%%%%%%%%%%%%
%%%%%%%%%%%%%%%%%%%%%%%%%%%%%%%%%%%%%%%%%%%%%%%%%%%%%%%%%%%%%%%%%%%%%%
\section{Introduction \& Motivation}
\label{sec:intro}

%%%%%%%%%%%%%%%%%%%%%%%%%%%%%%%%%%%%%%%%%%%%%%%%%%%%%%%%%%%%%%%%%%%%%%
\begin{frame} 
  \frametitle{Introduction}
  \begin{columns}
    \column{0.55\textwidth}
    \begin{itemize}
    \item Hard X-rays are strongly penetrating electromagnetic waves
    \item Interaction of X-ray and matter described by refractive
index
    \begin{equation*}
      n = 1 - \delta + \ii \beta
    \end{equation*} 
    {\small
    $\beta$: absorption, destructive\\ 
    $\delta$: soft-tissue sensitive
    }
  \item Use monochromatic coherent X-ray with energies
s.t.
    \begin{equation*} 
%      \frac{\delta}{\beta} \gg 1
      \delta / \beta \gg 1
    \end{equation*}
  \end{itemize}
  \column{0.45\textwidth}
  \small
  % Interaction
  \includegraphics[width=\textwidth]{figures/interaction}
  \vfill
  % delta over beta
   \psfragfig[width=0.5\textwidth,height=0.4\textheight]
    {figures/refractive-index/fig-ratio}
    { %\footnotesize
      \psfrag{Energy}[c][c]{\scriptsize Energy $E$ [\si{eV}]}
      \psfrag{ratio}[c][r]{\scriptsize Ratio $\delta / \beta$}
      \psfrag{iron}{\tiny Fe}
      \psfrag{calcium}{\tiny Ca}
      \psfrag{silicon}{\tiny Si}
      \psfrag{carbon}{\tiny C}
      \psfrag{waterwater}{\tiny H$_2$O}
    }
%  \includegraphics[width=\textwidth]{figures/beta_over_delta}
  \end{columns}
  \arrowright Imaging of weakly absorbing materials using phase-contrast 
\end{frame}

%%%%%%%%%%%%%%%%%%%%%%%%%%%%%%%%%%%%%%%%%%%%%%%%%%%%%%%%%%%%%%%%%%%%%%
\begin{frame}
  \frametitle{Motivation}
  \begin{itemize}
  \item Biological specimen w/o cartilage, bone, \& blood vessels are
    essentially pure-phase objects for $E>\SI{20}{keV}$
  \item Imaging techniques to study the development of model organisms
    are essential tools in Biology, Medicine, Toxicology, …
    \item Examples of important vertebrate model organisms
      \begin{tikzpicture}[]
        \def\xx{0.33\textwidth}
        \def\imh{0.22\textheight}
        \small
         \node[inner sep=1pt,above] at (-\xx,0){
           \includegraphics[height=\imh]{figures/model-organisms/xenopus}};         
         \node[inner sep=1pt,below] at (-\xx,0){Xenopus laevis};
         \node[inner sep=1pt,above] at (0,0){
           \includegraphics[height=\imh]{figures/model-organisms/zebrafish}};
         \node[inner sep=1pt,below] at (0,0){zebrafish};
         \node[inner sep=1pt,above] at (\xx,0){
           \includegraphics[height=\imh]{figures/model-organisms/lamprey}};
         \node[inner sep=1pt,below] at (\xx,0){lamprey};
    \end{tikzpicture}
  \end{itemize}
\end{frame}

%%%%%%%%%%%%%%%%%%%%%%%%%%%%%%%%%%%%%%%%%%%%%%%%%%%%%%%%%%%%%%%%%%%%%%
\begin{frame}
  \frametitle{Comparison: 4D imaging techniques}
  \begin{columns}[T]
    %% RIGHT
    \column{0.6\textwidth}
    \begin{itemize}
    \item Optical microscopy
      \begin{itemize}
      \item Non-destructive: Surfaces / thin samples
      \item Destructive: Sectioning
      \end{itemize}
    \end{itemize}
      \vfill
    \begin{itemize}
    \item Fluorescence microscopy
      \begin{itemize}
      \item Expression of fluorescence marker
      \item Photobleaching
      \item Sparse information
      \item Limited depth information: few cell layers
      \end{itemize}
    \end{itemize}
      \vfill
    \begin{itemize}
    \item Magnetic Resonance Imaging
      \begin{itemize}
      \item Non-destructive 
      \item Dense \& deep information
      \item Limited spatial \& temporal resolution
      \end{itemize}
    \end{itemize}
    %% LEFT
    \column{0.4\textwidth}
    \begin{tikzpicture}
    \def\yy{0.295\textheight}
    \def\imw{0.285\textheight}
    \small
    % top
    \node[left] at (0,2*\yy){
      \href{run:videos/comparison/xeno-optical-sectioning.avi?autostart&loop}{
      \includegraphics[width=\imw]{videos/comparison/xeno-optical-sectioning.jpg}}
    };         
    \node[right,align=left] at (0,2*\yy){Xenopus:\\ optical,\\ sectioned};
    % middle
    \node[left] at (0,1*\yy){
      \href{run:videos/comparison/zebra-light-sheet.avi?autostart&loop}{
        \includegraphics[width=\imw]{videos/comparison/zebra-light-sheet.jpg}}
    };         
    \node[right,align=left] at (0,1*\yy)
    {Zebrafish: \\light-sheet \\microscopy};
    % bottom
    \node[left] at (0,0*\yy){
      \href{run:videos/comparison/xeno-mri.avi?autostart&loop}{
        \includegraphics[width=\imw]{videos/comparison/xeno-mri.jpg}}
    };         
    \node[right,align=left] at (0,0*\yy)
    {Xenopus: \\MRI};
    \end{tikzpicture}
  \end{columns}
\end{frame}

%%%%%%%%%%%%%%%%%%%%%%%%%%%%%%%%%%%%%%%%%%%%%%%%%%%%%%%%%%%%%%%%%%%%%%
\begin{frame}
  \frametitle{Comparison: 4D imaging techniques}
    \begin{itemize}
    \item X-ray phase-contrast tomography
      \begin{itemize}
      \item Fast
      \item Dense 3D information
      \item Apt for optically opaque objects
      \item Possible time lapse $\lesssim$ \SI{2}{h}
      \end{itemize}    
    \end{itemize}    
    \begin{columns}[T]
      \column{0.6\textwidth}
    \href{run:videos/comparison/nature-movie_03.avi?autostart&loop}{
      \includegraphics[width=\textwidth
      ]{videos/comparison/nature-movie_03.jpg}} 
    \column{0.4\textwidth}
     %Time-lapse sequence of 
    Xenopus gastrulation:
    {\footnotesize 
      2-BM@APS\\ 
      energy $E=\SI{30}{keV}$,
      object-detector dist. $z$=\SI{62}{cm}\\
      pixel size $\Delta x$=\SI{1}{\micro\metre}\\
      bandwidth $\Delta E/E$=\num{e-2}\\
      exposure time/image $\Delta t$=\SI{15}{ms}\\
      time lapse $\sim\SI{2}{h}$
    }
  \end{columns}
  \srefr{Moosmann et al., Nature;
    Moosmann et al., Nature Protocols;
    Nawy, Nature Methods (research highlight)}
\end{frame}

%%%%%%%%%%%%%%%%%%%%%%%%%%%%%%%%%%%%%%%%%%%%%%%%%%%%%%%%%%%%%%%%%%%%%%
\begin{frame}
  \frametitle{Digital study}
  Using dense \& deep 3D information allows to digitally study
  \begin{itemize}
  \item Structural changes, gene expression, perturbation experiments:
    signal transduction and functioning of gene regulatory networks
  \item Cell and tissue dynamics: exchange of fluids and nutrients,
    deeply buried transient structures, cell migratory behaviour, ...
  \item Propulsion mechanisms for cell and tissue migration:
    differential flow-field analysis (convergent extension, cell-cell
    adhesion, cell-matrix interaction, ...)
  \end{itemize}
\end{frame}

%%%%%%%%%%%%%%%%%%%%%%%%%%%%%%%%%%%%%%%%%%%%%%%%%%%%%%%%%%%%%%%%%%%%%%
%%%%%%%%%%%%%%%%%%%%%%%%%%%%%%%%%%%%%%%%%%%%%%%%%%%%%%%%%%%%%%%%%%%%%%
\section{Experimental setup \& principles of 
propagation-based phase contrast}
\label{sec:setup-and-principles}

%%%%%%%%%%%%%%%%%%%%%%%%%%%%%%%%%%%%%%%%%%%%%%%%%%%%%%%%%%%%%%%%%%%%%%
\begin{frame}
  \frametitle{Experimental setup}
      \includegraphics[width=1\textwidth]{figures/setup/setup.jpg}
\end{frame}

%%%%%%%%%%%%%%%%%%%%%%%%%%%%%%%%%%%%%%%%%%%%%%%%%%%%%%%%%%%%%%%%%%%%%%
\begin{frame}
  \frametitle{Propagation-based phase contrast}
  \begin{columns}
    \column{0.5\textwidth}
    \begin{itemize}
    \item Simple setup: no additional optics, no scanning, no stepping
      $\rightarrow$ stable, fast, efficient
    \item Intensity contrast emerges via free-space propagation of
      transmitted wave field
    \end{itemize}
    \column{0.5\textwidth}
      \includegraphics[
      width=\textwidth
      ]{figures/setup/propagation-based-imaging.png}
  \end{columns}
  \vspace{0.5cm}
  Sequence of intensity maps $I_z$: signal increases
  \begin{tikzpicture}[]
    \def\xx{0.25\textwidth};
    \def\picheight{0.18\textwidth};
    \small
    % 5 cm    
    \node[inner sep=1pt,above] at (0,0){
      \includegraphics[height=\picheight]{figures/IntensitySequenceANKA/int0005mm.png}};         
    \node[inner sep=1pt,below] at (0,0){$\SI{0.5}{cm}$};
    % 18 cm    
    \node[inner sep=1pt,above] at (\xx,0){
      \includegraphics[height=\picheight]{figures/IntensitySequenceANKA/int0180mm.png}};         
    \node[inner sep=1pt,below] at (\xx,0){$\SI{18}{cm}$};
    % 58 cm    
    \node[inner sep=1pt,above] at (2*\xx,0){
      \includegraphics[height=\picheight]{figures/IntensitySequenceANKA/int0580mm.png}};         
    \node[inner sep=1pt,below] at (2*\xx,0){$\SI{58}{cm}$};
    % 108 cm    
    \node[inner sep=1pt,above] at (3*\xx,0){
      \includegraphics[height=\picheight]{figures/IntensitySequenceANKA/int1080mm.png}};         
    \node[inner sep=1pt,below] at (3*\xx,0){$\SI{108}{cm}$};
    \draw[very thick,->] (-0.5*\xx,-0.4) 
    node[below right]{object-detector distance $z$} 
    -- (3.5*\xx,-0.4);
    \end{tikzpicture}
\end{frame}

%%%%%%%%%%%%%%%%%%%%%%%%%%%%%%%%%%%%%%%%%%%%%%%%%%%%%%%%%%%%%%%%%%%%%%
%%%%%%%%%%%%%%%%%%%%%%%%%%%%%%%%%%%%%%%%%%%%%%%%%%%%%%%%%%%%%%%%%%%%%%
\section{Phase retrieval}
\label{sec:phase-retrieval}

%%%%%%%%%%%%%%%%%%%%%%%%%%%%%%%%%%%%%%%%%%%%%%%%%%%%%%%%%%%%%%%%%%%%%%
\begin{frame}
  \frametitle{Phase contrast}
  \begin{itemize}
  \item Rewrite Fresnel diffraction integral \sref{Guigay 1977}
    \begin{equation*}
    \F I_z(\vecxi) = 
    \int \Intd{\vecx}\exp{(-\ii 2\pi\vecx\cdot\vecxi)} 
    \Psi_0\left(\vecx-\tfrac{\lambda z}{2}\vecxi\right) 
    \Psi^*_0\left(\vecx+\tfrac{\lambda z}{2}\vecxi\right)
  \end{equation*}
  \item Expand exponential of exit wave field
    $\Psi_0=\sqrt{I_0}\expp{\ii\phi_0}$
    %=\sqrt{I_0}(1+\ii\phi_0-\phi_0^2+\order(\phi_0^3))
  \item Guigay up to $\order(\phi_0)$ (pure-phase object)
    \begin{equation*}
      \begin{split}
        &\F g_z(\vecxi) =
        2\sin(\pi\lambda z\vecxi^2)\F \phi_0(\vecxi) 
        \quad \leftarrow \quad \mathrm{linear\;order}
        \\ & -\cos(\pi\lambda z\vecxi^2)\int\Intd{\vecxi'}
        \F\phi_0(\vecxi')\,\F\phi_0(\vecxi-\vecxi')
        \\ & +\expp{\ii\pi\lambda z\vecxi^2}\int\Intd{\vecxi'}
        \expp{-\ii 2\pi\lambda z\vecxi\cdot\vecxi'} 
        \F\phi_0(\vecxi')\,\F\phi_0(\vecxi-\vecxi')
%        \\ & +\order( ( \ft{\phi}_0 )^3 ) \;.
      \end{split}
    \end{equation*}
  \end{itemize}
  \flr{intensity contrast $g_z=I_z/I_0=1$; $\vecx,\vecxi \in \R^2$}
\end{frame}

%%%%%%%%%%%%%%%%%%%%%%%%%%%%%%%%%%%%%%%%%%%%%%%%%%%%%%%%%%%%%%%%%%%%%%
\begin{frame}
  \frametitle{Phase retrieval}
  Linear models
  \vfill
  \begin{itemize}
  \item Contrast transfer function (CTF)
    \begin{itemize}
    \item Small relative phase variations:
      $\abs{\phi_0(\vecx+\tfrac{\lambda z}{2}\vecxi)
        -\phi_0(\vecx-\tfrac{\lambda z}{2}\vecxi)}\ll 1$
    \item High resolution
    \end{itemize}
  \end{itemize}

  \begin{equation*}
    \F\phi_0(\vecxi) = \frac{\F g_z(\vecxi) }{2\sin(\pi\lambda z\vecxi^2)}
  \end{equation*}
  \srefr{Guigay 1977; Gureyev et al. 2006; Cloetens 1998}
  \vfill
  \begin{itemize}
  \item Paganin phase retrieval
    \begin{itemize}
    \item Edge-enhancement regime: \\small propagation distance
    \item Low-pass filter
    \item Robust
    \end{itemize}
    \begin{equation*}
      \F\phi_0(\vecxi) = \frac{\F g_z(\vecxi) }{2\pi\lambda z\vecxi^2}
    \end{equation*}
    \srefr{Teague 1982; Paganin \& Nugent 1998; Bronnikov 2003}
  \end{itemize}
\end{frame}

%%%%%%%%%%%%%%%%%%%%%%%%%%%%%%%%%%%%%%%%%%%%%%%%%%%%%%%%%%%%%%%%%%%%%%
\begin{frame}
  \frametitle{Evolution of contrast transfer}
  \begin{itemize}
  \item Compute intensity contrast from input phase map (Lena test pattern)
  \item Amplify phase variations by upscaling of input phase
  \item Consider Fourier transformed intensity, \ie{} CTF, as a
    function of distance
  \end{itemize}
  \begin{columns}
    %% LEFT: Movie
    \column{0.4\textwidth}
    \begin{tikzpicture}
      \def\imh{0.9\textwidth}
      \small
      \node[inner sep=1pt,above] at (0,0){
        \href{run:videos/intensity/FT-intensity_vs_distance.avi?autostart&loop}{
          \includegraphics[height=\imh]
          {videos/intensity/FT-intensity_vs_distance.jpg}}
      };         
      \node[inner sep=1pt,below] at (0,0){$\abs{\F I_z}(\vecxi)$};
    \end{tikzpicture}
    %% RIGHT
    \column{0.6\textwidth}
  \begin{tikzpicture}
    \def\xx{0.55\textwidth};
    \def\imh{0.45\textwidth};
    \small
    % Image    
    \node[inner sep=0pt] at (0,0){
      \includegraphics[height=\imh]{figures/minimum/angint-abs-ft-int}};         
    \node[inner sep=0pt,above,rotate=90] at (-\imh/2,0){$\int\Intd\theta \abs{\F I_z}$};
    \node[inner sep=1pt,below] at (0,-\imh/2){$\abs{\vecxi}/\xi_{\mathrm{c}}$};
    % Image    
    \node[inner sep=0pt] at (\xx,0){
      \includegraphics[height=\imh]{figures/minimum/minima-pos}};         
    \node[inner sep=0pt,above,rotate=90] at (\xx-\imh/2,0){min. pos. $\abs{\vecxi}_{\mathrm{min}}/\xi_{\mathrm{c}}$};
    \node[inner sep=1pt,below] at (\xx,-\imh/2){Upscaling $S$};
    \end{tikzpicture}
  \end{columns}
\end{frame}

%%%%%%%%%%%%%%%%%%%%%%%%%%%%%%%%%%%%%%%%%%%%%%%%%%%%%%%%%%%%%%%%%%%%%%
\begin{frame}
  \frametitle{Phase retrieval: quasiparticle approach}
  \begin{columns}[T]
    \column{0.35\textwidth}
    \begin{itemize}
    \item Fourier space filter: crop regions around minima
    \item Phase retrieval via CTF using filtered intensity
    \end{itemize}

    \column{0.65\textwidth}
  \begin{tikzpicture}[]
    \def\xx{0.55\textwidth};
    \def\imh{0.37\textwidth};
    \def\yy{1.3*\imh}
    \small
    % intensity
    \node[inner sep=1pt,above] at (\xx,0){
      \includegraphics[height=\imh]{figures/ctf-filtering/CTFfiltering_intensity}};         
    \node[inner sep=1pt,below] at (\xx,0){$I_z$};
    % FT intensity
    \node[inner sep=1pt,above] at (2*\xx,0){
      \includegraphics[height=\imh]{figures/ctf-filtering/CTFfiltering_intensityFTabs}};         
    \node[inner sep=1pt,below] at (2*\xx,0){$\abs{\F I_z}$};
    % binary filter
    \node[inner sep=1pt,above] at (\xx,-\yy){
      \includegraphics[height=\imh]{figures/ctf-filtering/CTFfiltering_binaryfilter}};         
    \node[inner sep=1pt,below] at (\xx,-\yy){binary filter};
    % FT intensity
    \node[inner sep=1pt,above] at (2*\xx,-\yy){
      \includegraphics[height=\imh]{figures/ctf-filtering/CTFfiltering_intensityFTabsBinFilt}};         
    \node[inner sep=1pt,below] at (2*\xx,-\yy){filtered $\abs{\F I_z}$};
    \end{tikzpicture}
    \begin{flushright}
      {\scriptsize Data: TopoTomo@ANKA}
    \end{flushright}
  \end{columns}
\end{frame}

%%%%%%%%%%%%%%%%%%%%%%%%%%%%%%%%%%%%%%%%%%%%%%%%%%%%%%%%%%%%%%%%%%%%%%
\begin{frame}
  \frametitle{Tomography: Paganin vs CTF vs QP} 

  Xenopus, {\small4-cell stage, ID19@ESRF, $E=\SI{20}{keV}$, $\Delta
    E/E=\num{e-4}$, $z=$\SI{1}{m}, $\Delta
    x=\SI{0.75}{\micro\metre}$}
  
\begin{tikzpicture}[]
    \def\imh{0.28\textwidth}
    \small
    \node[inner sep=1pt,above] at (0,0){
      \includegraphics[height=\imh]{figures/paganin-ctf-qp/rend-pag-qp}};         
    \node[inner sep=1pt,below] at (-0.9*\imh,0){Rendering};
    \node[inner sep=1pt,below] at (0.1*\imh,0){Paganin};
    \node[inner sep=1pt,below] at (\imh,0){Quasiparticle};
    \end{tikzpicture}

  \begin{tikzpicture}[]
    \def\imh{0.3\textwidth}
    \small
    \node[inner sep=1pt,above] at (0,0){
      \includegraphics[height=\imh]{figures/paganin-ctf-qp/pag-ctf-qp}};         
    \node[inner sep=1pt,below] at (-\imh,0){Paganin};
    \node[inner sep=1pt,below] at (0,0){CTF};
    \node[inner sep=1pt,below] at (\imh,0){Quasiparticle};
    \end{tikzpicture}
\end{frame}


%%%%%%%%%%%%%%%%%%%%%%%%%%%%%%%%%%%%%%%%%%%%%%%%%%%%%%%%%%%%%%%%%%%%%%
%%%%%%%%%%%%%%%%%%%%%%%%%%%%%%%%%%%%%%%%%%%%%%%%%%%%%%%%%%%%%%%%%%%%%%
\section{In vivo imaging}
\label{sec:in-vivo}

%%%%%%%%%%%%%%%%%%%%%%%%%%%%%%%%%%%%%%%%%%%%%%%%%%%%%%%%%%%%%%%%%%%%%%
\begin{frame}
  \frametitle{In vivo imaging: Considerations}
  \begin{itemize}

  \item Motion blurring due to cell movement: $\sim \SI{4}{\micro\metre}$
    \begin{itemize}
    \item Acquisition time for tomogram at  resolution: $\lesssim \SI{30}{s}$
    \item Required flux density: $> \SI{e12}{photons/s/mm^2}$
    \end{itemize}

  \item Dose effects: heat load, double-strand breaking, free radicals
    induced by radiolysis of water

  \item Setup optimisation
    \begin{itemize}
    \item Energy: dose vs photon flux vs phase contrast vs detection efficiency
    \item Exposure time: dose vs signal-to-noise vs motion blurring
    \item Object-detector distance: signal-to-noise vs blurring
      (finite source size, coherence, nonlinear propagation effects)
    \item Tomography: number of projection, fast shutter influencing
      dose, resolution, rotation blurring, ...
    \end{itemize}

  \end{itemize}
\end{frame}

%%%%%%%%%%%%%%%%%%%%%%%%%%%%%%%%%%%%%%%%%%%%%%%%%%%%%%%%%%%%%%%%%%%%%%
\begin{frame}
  \frametitle{4D in vivo tomography} 
  \begin{columns}
    \column{0.45\textwidth}
    \vfill
    Xenopus gastrulation:
    {\small 2-BM@APS, time lapse: \SI{2}{h}, time
    between tomograms: \SI{10}{min}, $\Delta x$=\SI{1}{\micro\metre},
    $E$=\SI{30}{keV}, $\Delta E/E$=\num{e-2}}
    \vfill
  \begin{tikzpicture}[]
    \def\xx{0.55\textwidth};
    \def\yy{0.61\textwidth};
    \def\imh{0.49\textwidth};
    \small
        % raw proj
         \node[inner sep=1pt,above] at (\xx,0){
           \includegraphics[height=\imh]{figures/xeno-time-lapse/proj-raw}};         
         \node[inner sep=1pt,below] at (\xx,0){raw projection};
         % reco int
         \node[inner sep=1pt,above] at (0,-\yy){
           \includegraphics[height=\imh]{figures/xeno-time-lapse/reco-int}};         
         \node[inner sep=1pt,below] at (0,-\yy){intensity};
         % reco int
         \node[inner sep=1pt,above] at (\xx,-\yy){
           \includegraphics[height=\imh]{figures/xeno-time-lapse/reco-phase}};         
         \node[inner sep=1pt,below] at (\xx,-\yy){phase};
         \node[inner sep=11pt,below] at (\xx/2,-\yy){tomographic reconstruction};
    \end{tikzpicture}

    \column{0.55\textwidth}
    \begin{tikzpicture}
      % Animated rendering
      \node[inner sep=1pt,above] at (0,0){
        \href{run:videos/xeno-time-lapse/Nature_TomoSlicing.avi?autostart&loop}{
          \includegraphics[width=\textwidth
          ]{videos/xeno-time-lapse/Nature_TomoSlicing.jpg}}};         
      \node[inner sep=1pt,below] at (0,0){Slicing through rendered volume};
      % time-lapse sequence
      \def\yy{0.45\textheight}
      \node[inner sep=1pt,above] at (0,-\yy){
        \href{run:videos/xeno-time-lapse/Nature_TimeLapseSlice.avi?autostart&loop}{
          \includegraphics[width=\textwidth
          ]{videos/xeno-time-lapse/Nature_TimeLapseSlice.jpg}}};         
      \node[inner sep=1pt,below] at (0,-\yy){Time-lapse of reconstructed central slice};
  \end{tikzpicture}
  \end{columns}
\end{frame}

%%%%%%%%%%%%%%%%%%%%%%%%%%%%%%%%%%%%%%%%%%%%%%%%%%%%%%%%%%%%%%%%%%%%%%
\begin{frame}
    \def\xx{0.5\textwidth}
    \def\imh{0.35\textwidth}
    \def\imw{0.8\textwidth}
  \frametitle{New observation: transient structure}
  \begin{itemize}
  \item Confrontation of head \& ventral mesendoderm: formation of a
    transient ridge
    \item Not observable in fixed embryos or implants
  \end{itemize}
  \vfill
  \begin{tikzpicture}[]
    \small
    % im sequence
    \node[inner sep=1pt,above] at (0,0){
      \includegraphics[width=\imw]{figures/transient-ridge/trans-struct-sequ}};         
    \node[inner sep=1pt,below] at (0,0){Time lapse of reconstructed slice};
  \end{tikzpicture}
  \vfill
  \begin{tikzpicture}
    \small
    % embryo cut
    \node[inner sep=1pt,above] at (0,0){
      \includegraphics[height=\imh]{figures/transient-ridge/trans-struct-embryo-cut}};
    \node[inner sep=1pt,below,align=center] at (0,0){Moving contour of \\leading-edge mesendoderm};
    % segmentation
    \node[inner sep=1pt,above] at (\xx,0){
      \href{run:videos/transient-ridge/trans-struct-segm.avi?autostart&loop}{
        \includegraphics[height=\imh]{videos/transient-ridge/trans-struct-segm.jpg}}};
    \node[inner sep=1pt,below] at (\xx,0){Segmented structures};
  \end{tikzpicture}
\end{frame}

%%%%%%%%%%%%%%%%%%%%%%%%%%%%%%%%%%%%%%%%%%%%%%%%%%%%%%%%%%%%%%%%%%%%%%
\begin{frame}
  \frametitle{Data analysis}
  \begin{columns}
    %% RIGHT
    \column{0.5\textwidth}
    \begin{itemize}
    \item Visual inspection
    \end{itemize}
    \begin{tikzpicture}
      \def\imw{1.1\textwidth}
      \node[inner sep=1pt,above] at (0,0){
        \includegraphics[width=\imw]{figures/analysis-overview/arch-form}};
      \node[inner sep=1pt,below] at (0,0){
        \includegraphics[width=\imw]{figures/analysis-overview/trans-ridge}};
    \end{tikzpicture}
    \begin{itemize}
    \item Segmentation
    \end{itemize}
      \begin{tikzpicture}
        \def\imh{0.3\textheight}
        \node[inner sep=1pt,left] at (0,0){
          \includegraphics[height=\imh]{figures/analysis-overview/seg-cav}};
        \node[inner sep=1pt,right] at (0,0){
          \includegraphics[height=\imh]{figures/analysis-overview/seg-neurula}};
      \end{tikzpicture}
    %% MIDDLE
    \column{0.25\textwidth}
    % top
        \vfill
    \begin{itemize}
    \item Optical flow: overall dynamics
    \end{itemize}        
    \vfill
    \begin{itemize}
    \item Differential flow: collective vs individual movement
    \end{itemize}
    \vfill
    \begin{itemize}
    \item Integrated flow:\\ cell tracking
    \end{itemize}        
    %% RIGHT
    \column{0.25\textwidth}
    \def\imw{\textwidth}
    \includegraphics[width=\imw]{figures/analysis-overview/flow}
    \vfill
    \includegraphics[width=\imw]{figures/analysis-overview/flow-diff}
    \vfill
    \includegraphics[width=\imw]{figures/analysis-overview/flow-int}

    %   \begin{itemize}
    
    % \end{itemize}
    %   \begin{tikzpicture}
    %     \def\imw{0.5\textwidth}
    %     \node[inner sep=1pt,above] at (0,0){

    %   \end{tikzpicture}
  \end{columns}
\end{frame}

%%%%%%%%%%%%%%%%%%%%%%%%%%%%%%%%%%%%%%%%%%%%%%%%%%%%%%%%%%%%%%%%%%%%%%
\begin{frame}
  \frametitle{Segmentation \& volume analysis}
  \begin{tikzpicture}[]
    \def\xx{0.2\textwidth}
    \def\imh{0.3\textheight}
    \small
    \node[inner sep=1pt,above] at (0,0){
      \includegraphics[height=\imh]{figures/archenteron-inflation/arch-infl}};         
    \node[inner sep=0pt,below] at (-0.63*\xx,0){Archenteron inflation};
    \node[inner sep=0pt,below] at (1.33*\xx,0){Volume balancing};
  \end{tikzpicture}
  \vfill
  \begin{tikzpicture}
    \def\imh{0.3\textheight}
    \def\xx{0.5\textwidth}
    \small
    % rotation
    \node[inner sep=1pt,above] at (0,0){
      \href{run:videos/archenteron-inflation/rend-cavities.avi?autostart&loop}{
        \includegraphics[height=\imh]{videos/archenteron-inflation/rend-cavities}}};
    \node[inner sep=1pt,below] at (0,0){Segmented cavities};
    % inflation
    \node[inner sep=1pt,above] at (\xx,0){
      \href{run:videos/archenteron-inflation/infl-arch.avi?autostart&loop}{
        \includegraphics[height=\imh]{videos/archenteron-inflation/infl-arch}}};
    \node[inner sep=1pt,below] at (\xx,0){Inflating archenteron};
  \end{tikzpicture}
  \vfill
  \begin{itemize}
  \item Conclusion: Archenteron inflation is driven by uptake of
    external water
  \end{itemize}
\end{frame}

%%%%%%%%%%%%%%%%%%%%%%%%%%%%%%%%%%%%%%%%%%%%%%%%%%%%%%%%%%%%%%%%%%%%%%
\begin{frame}
  \frametitle{Optical flow}
  \begin{itemize}
  \item Visualising overall dynamics
  \end{itemize}
  \vfill
  \begin{tikzpicture}[]
    \def\imw{0.8\textwidth}
    \def\xx{0.4\textwidth}
    \def\yy{0.3\textheight}
    \small
    % rend
    \node[] at (0,0){
      \includegraphics[width=\imw]{figures/optical-flow/rend}};         
    \node[align=right,anchor=east] at (-\xx,0){Rendering of\\ halved embryo};
    % flow
    \node[] at (0,-\yy){
      \includegraphics[width=\imw]{figures/optical-flow/flow}};         
    \node[align=right,anchor=east] at (-\xx,-\yy){Velocity field\\on central slab};
  \end{tikzpicture}
  \vfill
\end{frame}

%%%%%%%%%%%%%%%%%%%%%%%%%%%%%%%%%%%%%%%%%%%%%%%%%%%%%%%%%%%%%%%%%%%%%%
\begin{frame}
  \frametitle{Differential optical flow}
  \begin{itemize}
  \item Collective vs individual movement
  \item Propulsion mechanisms
  \end{itemize}
  \vfill
  \begin{tikzpicture}[]
    \def\imh{0.4\textheight}
    \def\xx{0.3\textheight}
    \small
    % flow
    \node[inner sep=1pt,above] at (-\xx,0){
      \includegraphics[height=\imh]{figures/optical-flow-diff/flow}};         
    \node[inner sep=1pt,below,align=center] at (-\xx,0){Magnitude of \\flow field};
    % diff flow
    \node[inner sep=1pt,above] at (\xx,0){
      \includegraphics[height=\imh]{figures/optical-flow-diff/diff-flow}};         
    \node[inner sep=1pt,below,align=center] at (\xx,0){Magnitude of \\differential flow field};
  \end{tikzpicture}
  \vfill
\end{frame}

%%%%%%%%%%%%%%%%%%%%%%%%%%%%%%%%%%%%%%%%%%%%%%%%%%%%%%%%%%%%%%%%%%%%%%
\begin{frame}
  \frametitle{Integrated optical flow}
  \begin{itemize}
  \item Cell tracking via time integration
  \end{itemize}
  \vfill
  \begin{tikzpicture}[]
    \def\imh{0.4\textheight}
    \small
    \node[inner sep=1pt,above] at (0,0){
      \includegraphics[height=\imh]{figures/optical-flow-int/cell-track}};         
    \node[inner sep=2pt,below] at (0,0){Trajectories of cell pairs};
  \end{tikzpicture}
  \vfill
\end{frame}

%%%%%%%%%%%%%%%%%%%%%%%%%%%%%%%%%%%%%%%%%%%%%%%%%%%%%%%%%%%%%%%%%%%%%%
\begin{frame}
  \begin{center}
    \Large  Thank you for your attention!    
  \end{center}
\end{frame}

%%%%%%%%%%%%%%%%%%%%%%%%%%%%%%%%%%%%%%%%%%%%%%%%%%%%%%%%%%%%%%%%%%%%%%
\begin{frame}
  \frametitle{Outlook}
  \begin{itemize}
  \item 4th generation synchrotron upgrades implementing multi-bend
    achromat lattices: drastically enhanced coherence $\rightarrow$
    enhanced contrast, allows larger $z$
  \item Cone-beam magnification using Kirkpatrick Baez mirrors, ...
  \item Optimising detector system: scintillator, microscope, camera,
    blurring, deconvolution, point spread function, ...
  \item Tomography: fast shutter, dose, noise, stepping vs on-the-fly,
    alternative reconstruction techniques: algebraic, iterative,
    neural networks, sparsity, dictionary learning, ...
  \item Combination of X-ray phase-contrast \& confocal fluorescence
    microscopy
  \item Systematic dose studies: dose delay effects, use of radical
    scavengers, find phenomenological function describing dose expression
    effects, ...
  \item Upcoming experiments: emigration \& delamination processes of
    neural crest cells as a model for on-set of cancer metastasis
  \item Perturbation experiments (wildtype vs morphant embyros): role
    of adhesion molecules, signalling pathways, transcription factors,
    gene regulatory networks, ...
  \end{itemize}
\end{frame}

%%%%%%%%%%%%%%%%%%%%%%%%%%%%%%%%%%%%%%%%%%%%%%%%%%%%%%%%%%%%%%%%%%%%%%
%%%%%%%%%%%%%%%%%%%%%%%%%%%%%%%%%%%%%%%%%%%%%%%%%%%%%%%%%%%%%%%%%%%%%%    
\end{document}

% test
